\section{Background}

\begin{algorithm}
\caption{Optimal Ate pairing over Barreto-Naehrig curves~\cite{hsp10}}
\begin{algorithmic}[1]
\Require{$P \in \G_1, Q \in \G_2$}
\Ensure{$a_{opt}(Q, P)$}
\State{Write $s = 6t + 2 = \sum_{i=0}^{L-1} s_i 2^i$ where $s_i \in \{-1, 0, 1\}$ and $s_L = 1$}
\State{$T \leftarrow Q, f \leftarrow 1$}
\For{$i = L - 2$ to $0$}
    \State{$f \leftarrow f^2 \cdot l_{T,T}(P)$}
    \State{$T \leftarrow [2]T$}
    \If{$s_i = 1$}
        \State{$f \leftarrow f \cdot l_{T,Q}(P)$}
        \State{$T \leftarrow T + Q$}
    \EndIf
    \If{$s_i = -1$}
        \State{$f \leftarrow f \cdot l_{T,-Q}(P)$}
        \State{$T \leftarrow T - Q$}
    \EndIf
\EndFor
\State{$Q_1 \leftarrow \pi_p(Q), Q_2 \leftarrow \pi_p^2(Q)$}
\State{$f \leftarrow f \cdot l_{T,Q_1}(P), T \leftarrow T + Q_1$}
\State{$f \leftarrow f \cdot l_{T,-Q_2}(P), T \leftarrow T - Q_2$}
\State{$f \leftarrow f^{p^{12} - 1 / r}$}
\State{\Return $f$}
\end{algorithmic}
\end{algorithm}

\subsection{Pairing Computation}

An elliptic curve pairing is a non-degenerate bilinear map $e: \G_1 \times \G_2 \rightarrow \G_T$, where $\G_1$ and $\G_2$ are groups of points on elliptic curves over finite fields, and $\G_T$ is a multiplicative group in an extension field. For practical cryptographic applications, we require that this map be efficiently computable. We will focus on the optimal ate pairing~\cite{hsp10} over Barreto-Naehrig (BN) curves, which is widely used due to its efficiency. We will specifically use the Ethereum's BN254 curve, defined over a 254-bit prime field $\Fq$, with groups $\G_1, \G_2$ of order $r$ and target group $\G_T$ contained in the 12th extension field $\Fe{12}$.

The computation of an optimal ate pairing---the focus of our implementation---consists of two primary phases:

\begin{enumerate}
    \item \textbf{Miller loop}: The core iterative computation that evaluates line functions at pairs of points from $\G_1$ and $\G_2$, accumulating intermediate results to produce a value in $\Fe{12}$.
    
    \item \textbf{Final Exponentiation}: This step raises the Miller loop output to the power $(q^{12}-1)/r$, where $r$ is the order of the groups, ensuring the result lies in the correct subgroup $\G_T$.
\end{enumerate}

While both phases present optimization opportunities, the constraint-based verification of these computations introduces unique challenges and possibilities that differ significantly from native implementations.

\subsection{Schwartz-Zippel lemma} \label{sec:bg:sz}

The Schwartz-Zippel lemma (or more precisely, the Demillo-Lipton-Schwartz-Zippel lemma~\cite{DemLip78,Sch79,Zip79,Sch80}) provides a probabilistic method for testing polynomial identities. The lemma states that for a non-zero polynomial $P(x_1, x_2, \ldots, x_n)$ of total degree $d$ over a finite field $\F$, when evaluated at uniformly random field elements, the probability that $P$ evaluates to zero is at most $d/|\F|$.

This probabilistic bound becomes negligible for cryptographically sized fields where $d \ll |\F|$, making it practical to verify polynomial identities by evaluation instead of coefficient-wise operations.

\subsection{Constraint-Based Computation Models}

\subsubsection{Computational Model}

Our optimization targets constraint-based computation models where programs are represented as systems of algebraic constraints over finite fields. We formalize this model as follows:

\textbf{Definition (Constraint System):} A constraint system over a finite field $\Fq$ consists of:
\begin{itemize}
\item A set of variables $\mathbf{v} = (v_1, \ldots, v_m) \in \Fq^m$
\item A partition of variables into:
\begin{itemize}
\item \textit{Public inputs} $\mathbf{x} \in \Fq^n$ (known to verifier)
\item \textit{Private witnesses} $\mathbf{w} \in \Fq^{m-n}$ (known only to prover)
\end{itemize}
\item A set of polynomial constraints $\{C_i(\mathbf{v}) = 0\}_{i=1}^{\ell}$ where each $C_i$ is a polynomial over $\Fq$
\end{itemize}

A valid computation satisfies all constraints when variables are assigned appropriate values.

\textbf{Witnesses and Hints:} In constraint systems, complex operations (e.g., field division, polynomial quotients) can be computed efficiently using \textit{witnesses} or \textit{hints}---auxiliary values provided by the prover but verified by constraints. For instance, to verify $a/b = c$:
\begin{enumerate}
\item Prover provides $c$ as a witness
\item Verifier checks the constraint $b \cdot c = a$
\end{enumerate}

This approach transforms expensive operations into cheap verifications, fundamentally altering the cost model.

\textbf{Examples of Constraint Systems:}
\begin{itemize}
\item \textbf{R1CS} (Rank-1 Constraint Systems)\cite{r1cs13}: Constraints of the form $(A\mathbf{v}) \circ (B\mathbf{v}) = (C\mathbf{v})$ where $\circ$ is element-wise multiplication
\item \textbf{Plonkish}\cite{plnk19}: Gate-based constraints with custom gates and copy constraints
\item \textbf{Cairo/AIR}\cite{cairo}: Algebraic Intermediate Representation with polynomial constraints over execution traces
\end{itemize}

\subsubsection{Cost Model}

We measure computational cost in terms of \textit{constraints} (denoted by \textbf{C}), representing the fundamental unit of work in constraint systems. One constraint typically corresponds to one multiplication gate or one algebraic relation to verify.

\textbf{Basic Field Operations in $\Fq$:}
\begin{itemize}
\item \textbf{Addition/Subtraction:} 0C (constraints are linear, free to compose)
\item \textbf{Multiplication:} 1C (one constraint)
\item \textbf{Constant multiplication:} 0C (absorbed into constraint coefficients)
\item \textbf{Division/Inversion:} 2C (one multiplication constraint + one equality check)
\end{itemize}

\textbf{Extension Field Operations:} For $\Fe{k}$ elements represented as $k$ coefficients in $\Fq$:
\begin{itemize}
\item \textbf{Addition:} 0C (component-wise, linear)
\item \textbf{Multiplication:} Traditional approach requires $O(k^2)$ base field multiplications, but our polynomial verification reduces this significantly
\item \textbf{Evaluation at point $x$:} $k \cdot 1C = kC$ (Horner's method: $k-1$ multiplications + $k-1$ additions)
\end{itemize}

\textbf{Commitment Overhead:} Beyond constraint costs, we account for \textit{calldata} or \textit{commitment overhead}---the number of field elements the prover must commit to (via Fiat-Shamir hashing in non-interactive proofs). For blockchain deployments, this overhead directly impacts transaction costs and is often the dominant factor.

\subsubsection{Cost Model Comparison}

Table \ref{tab:cost_comparison} contrasts costs between native implementations and constraint systems:

\begin{table}[h]
\centering
\caption{Cost comparison: Native vs Constraint-based computation}
\label{tab:cost_comparison}
\begin{tabular}{|l|c|c|c|}
\hline
\textbf{Operation} & \textbf{Native Cost} & \textbf{Constraint Cost} & \textbf{Ratio} \\
\hline
$\Fq$ Multiplication & 1$\times$ & 1C & 1$\times$ \\
$\Fq$ Inversion & 100$\times$ & 2C & $\approx$50$\times$ cheaper \\
$\Fe{12}$ Multiplication & 50$\times$ & 54C (traditional) & Similar \\
$\Fe{12}$ Multiplication & 50$\times$ & 39C (our method) & Better \\
\hline
\end{tabular}
\end{table}

\textbf{Key Insight:} Operations that are traditionally expensive (inversions, divisions) become cheap in constraint systems when computed via witnesses. This inverts traditional optimization priorities: we can afford inversions to save multiplications, enabling affine coordinate systems and other techniques impractical in native implementations.

Recognizing and exploiting these unique computational properties is essential for unlocking hidden performance potential. Youssef El Housni's pioneering work \textbf{Pairings in Rank-1 Constraint Systems}~\cite{PR1CS22} represents, to our knowledge, the first systematic exploration of these optimization opportunities, providing a breakthrough in the practicality of recursive proving by thoughtfully leveraging the distinctive cost model inherent to constraint-based systems.